\documentclass[11pt,a4paper]{article}
% Common preamble for the hint sheets (\input from each file)
\usepackage[utf8]{inputenc}
\usepackage[T1]{fontenc}
\usepackage{lmodern}
\usepackage{amsmath,amssymb,amsthm,mathtools}
\usepackage[a4paper,margin=2.3cm]{geometry}
\usepackage{enumitem}
\usepackage{xcolor}
\usepackage[hidelinks]{hyperref}
\renewcommand{\leq}{\leqslant}
\renewcommand{\le}{\leqslant}
\renewcommand{\geq}{\geqslant}
\renewcommand{\ge}{\geqslant}
\setlength{\parindent}{0pt}
\setlength{\parskip}{0.35em}
% \hintstitle{Sheet title}{Class work / Homework}
\newcommand{\hintstitle}[2]{\begin{center}{\LARGE\bfseries Hints: #1\par}\vspace{0.5em}{\large #2\par}\end{center}
\noindent\textit{The hints to each problem are cumulative. Read one, try again, and open the next only if you are still stuck.}\par\medskip\hrule\medskip}
\newcommand{\problem}[1]{\par\bigskip{\large\bfseries Problem #1}\par\nopagebreak\smallskip}
\newcommand{\subpart}[1]{\par\smallskip\textbf{(#1)}\par\nopagebreak}
\newcommand{\hint}[1]{\par\smallskip\noindent\textcolor{black!70}{\textbf{Hint #1.}}\enspace}
\newcommand{\note}[1]{\par\smallskip\noindent\textit{Note.} #1\par}
\providecommand{\lt}{<}
\providecommand{\gt}{>}
\providecommand{\tightlist}{\setlength{\itemsep}{0pt}\setlength{\parskip}{0pt}}
\emergencystretch=3em

\begin{document}
\hintstitle{Number theory -- 8}{Homework}

\problem{1}
\hint{1} Use Vieta jumping (the Markov tree). From a solution $(x,y,z)$, the triple $(x,y,3xy-z)$ is again a solution. Start from $(1,1,1)$ and always replace the smallest coordinate.
\hint{2} $(1,1,1)\to(1,1,2)\to(1,2,5)\to(1,5,13)\to(5,13,194)\to(13,194,7561)\to(194,7561,4400489)$. The last triple has all entries $>50$.

\problem{2}
\hint{1} Look for $n$ with $n+1$ a perfect square: then $n+1=x^2+0^2$ and $n+2=x^2+1^2$, and only $n=x^2-1$ remains.
\hint{2} Make $x^2-1$ equal to $2y^2=y^2+y^2$. The Pell equation $x^2-2y^2=1$ has infinitely many solutions.

\problem{3}
\hint{1} The infinite continued fraction $x$ satisfies $x=2207-\frac1x$, so $x=\frac{2207+\sqrt{2207^2-4}}2$ (the larger root, since $x>1$).
\hint{2} Recognise $2207$ as a Lucas number: with $\varphi=\frac{1+\sqrt5}2$, $\varphi^{k}+\varphi^{-k}=L_k$, and $L_{16}=2207$. So $x=\varphi^{16}$, and the answer is $\varphi^2=\frac{3+\sqrt5}2$.

\problem{4}
\hint{1} Solve the recurrence: $a_n=\frac{(2+\sqrt3)^n+(2-\sqrt3)^n}2$, i.e.\ $a_n=T_n(2)$ for the Chebyshev polynomial $T_n$.
\hint{2} For odd $k$, $a_m\mid a_{mk}$, because $\cos(k\theta)$ is a polynomial in $\cos\theta$ with no constant term when $k$ is odd (equivalently $x^k+y^k$ is divisible by $x+y$). Use $2015=5\cdot403$.
\hint{3} $a_5=362=2\cdot181$, so $181$ divides $a_{2015}$. ($a_{13}$ and $a_{31}$ give other valid answers.)

\problem{5}
\hint{1} The coefficients satisfy $a_{n+2}=2a_{n+1}+a_n$, $a_0=1$, $a_1=2$: they are the Pell numbers, $a_n=\frac{(1+\sqrt2)^{n+1}-(1-\sqrt2)^{n+1}}{2\sqrt2}$.
\hint{2} Prove $a_n^2+a_{n+1}^2=a_{2n+2}$, either from the closed form (using $(1+\sqrt2)(1-\sqrt2)=-1$) or by induction via the matrix $\begin{pmatrix}2&1\\1&0\end{pmatrix}^n$. So $m=2n+2$.

\problem{6}
\hint{1} $a_{n+1}$ is the larger root of $t^2-7a_nt+a_n^2+9=0$. This relation is symmetric in $t$ and $a_n$, so $a_{n-1}$ is the other (smaller) root.
\hint{2} Vieta: $a_{n+1}+a_{n-1}=7a_n$. With $a_0=1$, $a_1=5$, induction gives integrality (in fact $a_n=F_{4n+1}$).
\hint{3} From $a_{n+1}^2+a_n^2+9=7a_na_{n+1}$ you get $(a_n+a_{n+1})^2=9(a_na_{n+1}-1)$. Hence $a_na_{n+1}-1=\bigl(\frac{a_n+a_{n+1}}3\bigr)^2$, and $3\mid a_n+a_{n+1}$ follows from $9\mid(a_n+a_{n+1})^2$.

\problem{7}
\hint{1} If a prime $p$ divides two coordinates, it divides the third. Then it divides all three after any Vieta jump, so the property would survive the descent down to $(1,1,1)$. Hence the coordinates are pairwise coprime.
\hint{2} $z\mid x^2+y^2$ with $\gcd(x,y)=1$. So $4\nmid z$, since $x^2+y^2\equiv1,2\pmod4$, and $-1$ is a quadratic residue modulo every odd prime $p\mid z$, i.e.\ $p\equiv1\pmod4$.
\hint{3} For $x^2+y^2+z^2=kxyz$, take a solution with $x\le y\le z$ and minimal sum. The other root $z'=kxy-z$ must satisfy $z'\ge z$, hence $z^2\le zz'=x^2+y^2\le 2y^2$. Plug in: $kxy^2\le kxyz=x^2+y^2+z^2\le2(x^2+y^2)\le4y^2$, so $kx\le4$. Here $kx=4$ would force $x=y=z$ and $3x^2=4x^2$, so in fact $kx\le3$.
\hint{4} Cases: $k=3,x=1$ (Markov); $k=1,x\le3$; $k=2,x=1$, which gives $(y-z)^2=-1$, impossible. For $k=1$, look modulo $3$ at any solution: squares are $0,1$, and checking how many of $x,y,z$ are divisible by $3$ shows $3\mid x,y,z$. Then $(x/3,y/3,z/3)$ is a Markov triple; conversely $(3x,3y,3z)$ always works.

\problem{8}
\hint{1} Multiplying by $\pm\varepsilon^{k}$ moves any solution $\beta>0$ into $[\sqrt{|n|},\varepsilon\sqrt{|n|})$. Then $|\bar\beta|=|n|/\beta\le\sqrt{|n|}$. So $|u|=\frac{|\beta+\bar\beta|}2\le\frac{(\varepsilon+1)\sqrt{|n|}}2$ and $|v|\le\frac{(\varepsilon+1)\sqrt{|n|}}{2\sqrt d}$, leaving finitely many classes.
\hint{2} Counting: in each class the positive solutions are $\beta\varepsilon^k$ with $k$ large enough, and $x\approx\frac12\beta\varepsilon^k$. So $x\le X$ holds for $\log_\varepsilon X+O(1)$ values of $k$. Summing over the $h$ classes gives $\frac h{\ln\varepsilon}\ln X+O(1)$.
\hint{3} For $x^2-3y^2=13$ we have $\varepsilon=2+\sqrt3$. Search $\beta$ in the box from Hint 1: the classes are represented by $4+\sqrt3$ and $5+2\sqrt3$. All positive solutions are $(4+\sqrt3)(2+\sqrt3)^k$ and $(5+2\sqrt3)(2+\sqrt3)^k$, $k\ge0$: $(4,1),(5,2),(11,6),(16,9),(40,23),(59,34),\dots$

\problem{9}
\hint{1} Lagrange: if $\alpha$ is a quadratic irrational, its complete quotients have the form $\alpha_k=(P_k+\sqrt D)/Q_k$ with $Q_k\mid D-P_k^2$. Show that $P_k,Q_k$ are bounded, so some $\alpha_k$ repeats. The converse ($\alpha$ periodic implies quadratic) is the fixed point of a M\"obius map.
\hint{2} Galois: call $\alpha$ reduced if $\alpha>1$ and $-1<\bar\alpha<0$. Show that $\alpha$ reduced implies $\alpha_1$ reduced, and that a reduced number determines its predecessor ($-1/\bar\alpha_{k+1}$ has integer part $a_k$). So the period has no pre-period. Conjugation reverses the period.
\hint{3} $\sqrt d+a_0$ is reduced, which gives $\sqrt d=[a_0;\overline{a_1,\dots,a_{r-1},2a_0}]$; the palindrome comes from reversal. Compute: $\sqrt{m^2+1}=[m;\overline{2m}]$ with $\eta=m+\sqrt{m^2+1}$ (norm $-1$) and $\varepsilon=\eta^2$. Also $\sqrt{m^2-1}=[m-1;\overline{1,2m-2}]$ with $\eta=\varepsilon=m+\sqrt{m^2-1}$.

\problem{10}
\hint{1} Write $a(a+1)=Du^2$ with $D$ squarefree. Then $(a^2+a)(b^2+b)$ is a square if and only if $a$ and $b$ have the same $D$. Also $(2a+1)^2-4Du^2=1$, a Pell equation.
\hint{2} The diagonal $a=b$ gives $N$ pairs. An off-diagonal pair needs two solutions of $X^2-4DU^2=1$ with $X\le2N+1$. This forces the fundamental solution to be $\lesssim\sqrt N$, so only $O(\sqrt N)$ values of $D$ qualify, each with $O(\log N)$ solutions.
\hint{3} Hence $S_N=N+O(\sqrt N\log^2N)$ and $\ell=1$. For the second question: $b=4a(a+1)$ pairs with $a$, since $a(a+1)b(b+1)=\bigl(2a(a+1)(2a+1)\bigr)^2$. This gives $\gg\sqrt N$ off-diagonal pairs, so $|S_N-N|\le N^{1/2025}$ is \emph{false}.
\end{document}
